\chapter{Input  file. }

Input for  \TiberCAD is  composed by an  input  file e.g. "input.tib" and a mesh file generated by a mesher software: as for  now, mesh files from GMSH (*.msh,  v.1)  and from ISE-TCAD (*.grd) are  supported.



A valid  input file is a  text file with the  following  structure.
  
It is composed by several sections:
%
each \textbf{section} begins with a section-name preceded by \char`\"{}\$\char`\"{} (e.g. \$Physics ).
A section is enclosed between  \char`\"{}\{\char`\"{} and \char`\"{}\}\char`\"{} brackets and is composed 
 by a number of  blocks enclosed between  \char`\"{}\{\char`\"{} and \char`\"{}\}\char`\"{} brackets.
   Each  \textbf{block} can be possibly composed by one or more blocks, each preceded by a block-name. 
 Everything following a '\textbf{\#}' is considered as a comment and is disregarded.

A parameters-block is  an elementary block which can contain 
%Except the model-blocks and the BC\_\-regions-blocks (see in the following), each block can contain 
zero or any number of   parameter assignements in the  form:

 \char`\"{}\textit{tagname} = \textit{tagvalue}\char`\"{}, where \char`\"{}\textit{tagname}\char`\"{} is a string and \char`\"{}\textit{tagvalue}\char`\"{} is a single numerical or string item or a list of items between \char`\"{}(\char`\"{} and \char`\"{})\char`\"{} parenthesis. Format is free for these assignements, provided they are separated by spaces. Everything which follows a '\textbf{\#}' is considered as a comment and is disregarded.
 Here and in the whole input file a string item can include a combination of characters, special characters and numbers, but not spaces; if a space is found , the string item is taken as terminated.
 
  
%List  of  keywords:

The input file is composed by  the following sections:

 \textbf{Device}  ,   \textbf{Models}  ,  \textbf{Physics},  \textbf{Solver},   \textbf{Simulation} 

%\textit{Models}:  driftdiffusion , thermal  ,  excitontransport,  macrostrain,   sweep(special)

\section{Device section }
%\textbf{Device} section:
\begin{verbatim}
$Device
{
 Region buffer
   {
   .........
   }
 Region barrier_1
   {
     ............
   }
   .................
}


\end{verbatim}





  In \char`\"{}\textbf{Device}\char`\"{} section, each block (Region-block) must be preceded by the keyword \char`\"{}\textbf{Region}\char`\"{}, followed by the  (single-word) name of the physical region.
For an ISE  mesh it must  be  the  same name defined during the  modeling of the  device.

\begin{verbatim}
   Region QWell
   {
     reg_numb = 4
     structure = wz
     y-growth-direction = (1,0,-1,0)
     z-growth-direction = (-1,2,-1,0)
     x-growth-direction = (0,0,0,1)
     mat = GaN
     doping = 1e17 doping_type = donor doping_level = 0.025
   }
\end{verbatim}
%The following keywords must always be  present in the body of each Region-block:  mat, reg\_numb .
%Other keywords  are optionals.
Here are the  description of  the available keywords for a Region-block.

%In the  header of  each block:
%\textit{Region} :   name of the  region; for ISE grid mesh it should  be  the  name defined during the  modeling of the  device.

%In the  body of the  block:

\textit{mat} (mandatory) :  name  of the  material associated to the  present region; it  can be  an  alloy,  in this  case the  keyword \textit{x } must be  present.

\textit{x} : alloy concentration; e.g. in $Al_xGa_{1-x}As$,  x is  Al concentration

\textit{ reg\_numb} (mandatory) : region ID as  specified in the meshing  program (GMSH).

\textit{structure} : crystal structure ( wz = wurtzite,  zb =  zincblend)

\textit{x-growth-direction}, \textit{y-growth-direction}, \textit{z-growth-direction}: Bravais  vectors with Miller  indexes for  wurtzite crystal (4 element vectors) or zincblende crystal  (3 element vectors).

\textit{doping} : doping  concentration [cm$^{-3}$]

\textit{doping\_type} : donor or  acceptor

\textit{doping\_level}: energy level of  the dopant [eV]


\section{Models section }



\begin{verbatim}
$Models
{ 
 model driftdiffusion
  { 
   ..........
   BC_Regions  
    {
      BC_Region cathode
      {
       ............
      }
      BC_Region anode
      {
        ..............
      }
    }   
  } 
 model macrostrain
  { 
  ...........................
 
\end{verbatim}




   In \char`\"{}\textbf{Models}\char`\"{} section, one or more  model-blocks must be present: each model-block must be preceded by the keyword \char`\"{}\textbf{model}\char`\"{}, followed by the (single-word) model name. 
 This must be the   name of  one of the  \TiberCAD  simulation models.
 
Here are the simulation models implemented until now:  \textbf{driftdiffusion} , \textbf{thermal}  ,  \textbf{excitontransport},  \textbf{macrostrain}.
   
   Each model-block can  contain the following optional blocks:
   
 %  one or more optional Boundary Condition regions blocks (\textbf{BC\_\-regions}-block).
   
   
 %   \char`\"{}\textbf{Models}\char`\"{} section can contain also 
   
\begin{itemize}
	
   \item one    "`options"' block, preceded by the keyword \char`\"{}\textbf{options}\char`\"{}. This block can contain    general  options for the present  model.
	\item  one or more  \textbf{physical\_model} blocks: each physical\_model block must be preceded by the keyword \char`\"{}\textbf{physical\_\-model}\char`\"{}, followed by the (single-word) name of the physical model.
	Each physical\_model block can contain parameters  relevant to a specifical model of a physical property or quantity related to the present model.
\item    one or more  Boundary Condition regions blocks (\textbf{BC\_\-regions}-block).
   The BC\_\-regions-block must be preceded by the keyword \char`\"{}\textbf{BC\_\-Regions}\char`\"{} and it is composed by one or more parameters-blocks, each preceded by the keyword \char`\"{}\textbf{BC\_\-Region}\char`\"{} followed by the (single-word)  name of the boundary condition region. This parameters-block can contain the possible description of the boundary region.
\end{itemize}


 
A detailed description of the possible parameters for these blocks  follows.



\subsection{options block }

\begin{verbatim}
options
    {
     simulation_name = driftdiffusion
     physical_regions = all
    }
\end{verbatim}

\textit{simulation\_name}: user-defined name of the particular instance of  the simulation model defined for this block. More than one simulation (with different  name and properties) can be  defined,in separated model blocks, which refer to a same  \TiberCAD  simulation model. If simulation\_name is  not assigned,  by default the \TiberCAD model name is taken as current simulation\_name.

\textit{physical\_regions}: (list of) physical region(s) to  which the  present simulation model will be  applied.
Physical region(s) can be integer numbers, consistent with the physical regions set assigned through GMSH, or strings, corresponding to the  region names defined in ISE grid  file.
Default value  is  "`all"' (all  physical regions of the  device).



\subsection{physical\_model block }

\begin{verbatim}
    physical_model recombination
    {
      model = SRH
    }
\end{verbatim}


Currently available physical\_models and parameters:

\textbf{recombination}:

\textit{model}:  SRH,direct, exciton\_generation,  exciton\_dissociation

\textit{exciton\_simulation}: excitons


\textbf{electron\_mobility}

\textit{model}: doping\_dependent

\textbf{hole\_mobility}

\textit{model}: doping\_dependent



\subsection{BC\_region block }

\begin{verbatim}
 BC_Region anode
      {
        BC_reg_numb = 2
        type = ohmic
	      voltage = 0.0
      }
\end{verbatim}



%For each BC\_region:

\textit{BC\_region}: name of the  present boundary  region

\textit{ BC\_reg\_numb}:  BC region ID(s) as  specified in the meshing  program (GMSH).

\textit{ type}: type of  boundary condition: ohmic,  schottky,...  ,  substrate (for strain calculations).

\textit{ voltage}:  value of voltage [V] applied to the present BC region (for ohmic and  schottky BCs).

\textit{ zero\_grad\_fermi\_h}, \textit{ zero\_grad\_fermi\_e}: if true set  Neumann b.c. to the fermi level in the b.c. region.

If \textit{type} is \textsl{substrate}:

\textit{material} : name of  material in the \textsl{substrate}  region. 

\textit{structure} : crystal structure ( wz = wurtzite,  zb =  zincblend)

\textit{x-growth-direction}, \textit{y-growth-direction}, \textit{z-growth-direction}: Bravais  vectors with Miller  indexes for  wurtzite.




\section{Solver section }

\begin{verbatim}
$Solver
{ 
 driftdiffusion
  {
    quasi_equilibrium = (cathode, anode)
    nonlin_max_it = 25
    nonlin_abs_tol = 1e-7
  #  nonlin_step_tol = 1e-3
    lin_max_it = 150
    ksp_type = gmres
    pc_type = composite
    #pc_type = ilu
    ls_max_step = 0.05
    #current_integration_method = rstf
  }
...................

\end{verbatim}

In this section one can  define the  setting parameters for  the numerical solvers: the section is organized in blocks, each one preceded by a block-name equal to one of the  user-defined simulation name. Solver parameters defined in each block refer to that particular simulation.

Two extra blocks can be  present:
\begin{itemize}
	
	
%\begin{verbatim}
% BC_Region anode
%      {
%        BC_reg_numb = 2
%        type = ohmic
%	      voltage = 0.0
%      }
%\end{verbatim}
%
	
   \item \textbf{sweep} block, preceded by the keyword \char`\"{}\textbf{sweep}\char`\"{}. This block  contains  the  setting for a set of  calculations applied to a boundary region (e.g. to a drain contact for the calculation of a IV characteristic).
   
   
%\begin{verbatim}
% BC_Region anode
%      {
%        BC_reg_numb = 2
%        type = ohmic
%	      voltage = 0.0
%      }
%\end{verbatim}

  \item  \textbf{selfconsistent} block , preceded by the keyword \char`\"{}\textbf{selfconsistent }\char`\"{}. In this block it is possible to define a self-consistent calculation based on two different simulation models (e.g. driftdiffusion and excitontransport).

\end{itemize}

Here are some of  the  parameters:

\textbf{quasi\_equilibrium:}

\textbf{nonlin\_max\_it}

\textbf{nonlin\_abs\_tol}

\textbf{lin\_max\_it}

\textbf{ksp\_type}

\textbf{pc\_type}

\textbf{ls\_max\_step}

\textbf{ls\_type}

\textbf{current\_integration\_method}


%sweep
\textbf{simulation} :  simulations to which sweep is  applied

\textbf{boundary}: boundary region to which sweep is  applied

\textbf{start} , \textbf{stop}, \textbf{steps} : sweep starts from \textit{start} value, is repeated  \textit{steps} times and stops in \textit{stop}

\textbf{plotvariable}

%% selfconsistent

\textbf{simulations} : simulations to be   executed in a self-consistent cycle. 

%thermal
\textbf{current\_simulation}  :    current simulation model used for  thermal  model

%macrostrain
\textbf{substrate} 




\section{Physics section }
In this section several physical parameters can be entered,  in  addition to or  overwriting  the material parameters present in the material files.
The section is organized in blocks, each one preceded by a block-name equal to one of the  user-defined simulation name. Physical parameters defined in each block refer to that particular simulation.

Here are some of the possible parameters:

\textbf{statistics} : FD

\textbf{model}: specification of generical drift-diffusion or  exciton transport  simulation.
Possible values: \textit{simple}, \textit{strained} ,  \textit{unstrained}.  

\textbf{tau\_rad}, \textbf{tau\_nonrad},  tau\_diss

\textbf{R}

\textbf{eff\_mass}

\textbf{mobility}

\textbf{generation\_model}



\section{Simulation section }

In this  section one  can  specify  several general parameters and settings  for  the  actual  calculation to be  run, such as the mesh file to be  used, the dimension of  simulation, the process-flow of  simulation, etc.

\textbf{searchpath} :  path for  material  files 

\textbf{meshfile} : name of  mesh file. N.B.: the extension is mandatory! ( .grd for ISE-TCAD, .msh for gmsh v.1 )

\textbf{mesh\_units} : units of  measurements used in the  meshing (relative to meters): e.g., $10^{-6}$ for  $\mu m$

\textbf{dimension} :  dimension of  simulation (1,2,3)

\textbf{temperature} :  temperature  of the system

\textbf{solve} : list of  simulations to  be  executed, in the order of  execution; if the list   contains "`sweep"', a sweep is  performed as specified in \textbf{sweep} block in the  \textbf{Solver} section.   

\begin{verbatim}
  solve =  (strain,driftdiffusion, quantum_electrons, quantum_holes)
\end{verbatim}



\textbf{resultpath} :  path for output  directory 

\textbf{output\_format} : format of  the  output data: \textbf{gmv},  \textbf{tecplot}, \textbf{grace} (for xmgr, ascii data column type),  ....

\textbf{plot} :  list of output variables which are calculated and available in output  files.

